\section{Requerimientos del sistema de instrumentación}

Esta sección constituye el puente entre la geofísica y el diseño electrónico. Cada requerimiento que sigue se deriva de una necesidad del método expuesta en las secciones 5 a 7 y no de una preferencia de implementación.

\textbf{Banda útil de trabajo.} La profundidad de investigación se fija por la máxima longitud de onda observable a través de $z\approx\lambda/2$, y la longitud de onda se relaciona con la frecuencia mediante $\lambda=c/f$. Combinando ambas expresiones, la frecuencia mínima necesaria para alcanzar una profundidad $z$ resulta

\begin{equation}
f_{\min}\simeq\frac{c_R}{2z}.
\end{equation}

Con la velocidad de fase efectivamente observada en el sitio, del orden de 150 m/s, alcanzar 10 m exige contenido coherente en torno a 7,5 Hz, mientras que el objetivo formal de 50 m exigiría alcanzar aproximadamente 1,5 Hz. Esta única relación explica la distancia entre el objetivo de diseño y el alcance demostrado, y anticipa por qué la elección del transductor y la energía de la fuente, y no la velocidad de conversión, constituyen el cuello de botella del sistema.

% — verificar el criterio exacto contra la fuente antes del .bib
\textbf{Separación entre canales y apertura del arreglo.} La separación $\Delta x$ define el muestreo espacial del frente de onda y por lo tanto un límite de aliasing, ya que sólo pueden interpretarse sin ambigüedad longitudes de onda que satisfagan $\lambda\gtrsim 2\Delta x$. La apertura total $L$ actúa en el extremo opuesto: no puede estimarse de forma confiable una longitud de onda comparable o mayor que la extensión del tendido. Las recomendaciones de buena práctica expresan ese límite como $L\gtrsim 1{,}5\,\lambda_{\max}$ en una estimación optimista, y con criterios más conservadores exigen $L\gtrsim 2\,\lambda_{\max}$ \citep{Foti2018}. Para la apertura de 40 m efectivamente desplegada, el criterio optimista admite longitudes de onda de hasta unos 27 m, cifra que resulta consistente con el soporte de 20 a 22 m que el procesamiento identificó a posteriori y que se reporta en la sección 11.3. Ambos límites se trasladan al plano $f$–$c$ y deben dibujarse sobre la imagen de dispersión, práctica que se adopta en la sección 11.

\textbf{Sincronización entre canales.} Un desfase temporal $\Delta t$ entre dos canales introduce un error de fase $2\pi f\,\Delta t$ que el procesamiento interpreta como un tiempo de propagación adicional. En una aproximación de primer orden, el error relativo resultante sobre la velocidad de fase es

\begin{equation}
\frac{\Delta c}{c}\simeq\frac{c\,\Delta t}{\Delta x},
\end{equation}

de modo que la tolerancia de sincronización queda determinada por la geometría y por la velocidad esperada, no por una cifra elegida de antemano. Para $c\approx150$ m/s, $\Delta x=2$ m y un error admisible del 1 \% en la velocidad de fase, el desfase entre nodos debe mantenerse por debajo de unos 130 µs. Este número es el que justifica la arquitectura de arranque coordinado descrita en la sección 9 y el que deberá verificarse experimentalmente según se indica en la sección 12.

\textbf{Relación señal a ruido y rango dinámico.} El registro debe conservar simultáneamente el primer arribo cercano a la fuente y la llegada atenuada en el offset más lejano, cuya diferencia de amplitud abarca varios órdenes de magnitud. De allí la necesidad de ganancia programable por canal y de un convertidor con margen suficiente para evitar tanto la saturación como el enterramiento de la señal lejana en el ruido de cuantificación. La contrapartida es que ninguna ganancia recupera información que nunca superó el piso de ruido de la cadena, lo que devuelve la atención a la energía de la fuente.

\begin{figure}[!tb]
\centering
\includegraphics[width=\columnwidth,keepaspectratio]{sincronizacion_fase.png}
\caption{Traducción del desalineamiento temporal entre canales a error de fase y, a través de la geometría, a error de velocidad de fase. El presupuesto fija el criterio de medición de la sincronización.}
\end{figure}

\textbf{Preservación de fase.} Como la magnitud medida es una velocidad de fase, un error sistemático de fase entre canales se traduce directamente en un error de velocidad. El requerimiento no es que la respuesta del acondicionamiento sea plana, sino que sea conocida y, sobre todo, idéntica entre canales, porque una diferencia común a todos los canales se cancela al estimar la pendiente de fase mientras que una dispersión entre unidades no lo hace.

\textbf{Trazabilidad y operación en campo.} El procesamiento sólo es reproducible si cada registro conserva su geometría, su configuración de ganancia, su referencia temporal y la identificación del hardware que lo produjo. A ello se suma la exigencia operativa de trabajar sin infraestructura, con almacenamiento local y sin depender de conectividad externa durante la adquisición.

La tabla 1 resume esa derivación y funciona además como mapa de lectura del resto del documento: la columna de decisiones anticipa lo que la sección 9 desarrolla, y la de evidencia remite a las secciones 10 y 11, donde cada afirmación se sostiene o se declara pendiente.


\begin{table}[!ht]
\centering
\footnotesize
\caption{Puente requerimiento → decisión → evidencia → estado. Es la tabla central del documento y no debe recortarse.}
\begin{tabularx}{\columnwidth}{YYYY}
\toprule
Necesidad derivada del método & Decisión de diseño adoptada & Evidencia disponible & Estado \\
\midrule
Observar longitudes de onda largas ($f\lesssim 8$ Hz para 10 m) & Geófono SM-24 con rama de compensación $1-\mathrm{BP}$ \citep{Ma2023} → §9.3 & Ruta PGA → ADC identificada desde 0,21 Hz → §10.1 & Electrónica verificada en la banda; el geófono entra como modelo \\
Evitar aliasing espacial y respetar la apertura & Geometría registrada: $\Delta x=2$ m, apertura 40 m, offsets de 10 a 50 m & Metadatos de la campaña & Disponible para la campaña realizada \\
Estimar $c_R(f)$ sin sesgo instrumental & Respuesta de fase conocida y homogénea entre canales & Caracterización parcial de fase del AFE & Pendiente: falta dispersión entre placas \\
Sostener el rango dinámico entre offsets & Ganancia programable, control de offset y ADC $\Delta\Sigma$ con rangos seleccionables & Arquitectura implementada y registros de campo utilizables & Pendiente: sin ruido referido a entrada ni ENOB \\
Sincronizar por debajo de ~130 µs & Arranque coordinado por radio y temporización de ventana en lógica dedicada → §9.4 & Ensayo funcional con tres nodos en laboratorio → §10.2 & Funcionalmente demostrado; sin verificación metrológica \\
Referenciar temporalmente cada impacto & Nodo HAMMER con martillo instrumentado, adquirido junto al nodo GEO & 598 registros aceptados con referencia de impacto & Operativo en campaña \\
Adquirir simultáneamente varios receptores & Nodos GEO independientes con buffer local e inicio común → §9.1 & Coordinación probada con tres nodos en laboratorio → §10.2 & No demostrado con varios GEO completos en campo \\
Evitar pérdida o corrupción de muestras & Fragmentación con numeración de secuencia, reintento, CRC y copia local & Implementación e interfaces existentes & Falta ensayo multinodo bajo interferencia \\
Conservar trazabilidad del registro & Metadatos de campaña, manifiesto y servidor de ingesta & Conjunto de datos y flujo reproducible & Falta asociar revisión de hardware y firmware a cada registro \\
\bottomrule
\end{tabularx}
\end{table}
